\section{Conclusion}
\label{sec:conclusion}

The argument of this paper, in one sentence: hierarchical memory inside an industrial AI inference node is regime-dependent in a way that can be written down, and once you write it down the engineering problem becomes a lot more tractable.

Concretely, three regimes---capacity-limited, coordination-dominated, and I/O-limited---govern how a serving node behaves under industrial load. RAMSES, the orchestrator we built around that view, coordinates VRAM, DRAM, and NVMe under thresholds derived from the latency decomposition. The boundaries are not heuristics; they fall out of the model, and we then validated them through controlled experiments and long-running industrial traces.

Numbers from the runs: load time down by up to 60\%, inference latency down 35\%, energy per token down 16\%, SLA violations down by an order of magnitude, all under sustained industrial workload across heterogeneous models and multi-GPU configurations. Those gains live inside the analytically identifiable coordination-dominated region. Outside it, they are bounded---which is also a useful thing to know, because it tells a deployment engineer where coordination stops paying back.

The wider point is practical. Holding memory-induced jitter down keeps inference inside the TSN window. Inference inside the TSN window means fewer missed scan cycles. Fewer missed scan cycles, in turn, means less risk of production downtime, all while throughput-per-watt improves under the kind of tight industrial power budget that real factory floors actually run under. We treat this work as a starting point---a principled foundation, but not the final word---for deterministic, cost-efficient AI infrastructure inside industrial informatics.
